\chapter{Tests, Validation Expérimentale et Résultats}

\section{Stratégie Globale de Test et d'Assurance Qualité}

La fiabilité d'un système de supervision industrielle temps réel repose sur une rigueur de validation stricte. Une stratégie de test pyramidale a été mise en œuvre afin de garantir l'absence de régression et l'exactitude des algorithmes décisionnels :

\begin{itemize}
    \item \textbf{Tests Unitaires du Cœur Métier :} Vérification exhaustive des invariants de domaine, des mutations d'état et des exceptions de garde ;
    \item \textbf{Tests des Algorithmes de Graphe :} Validation formelle du parcours BFS sur diverses topologies industrielles (graphes linéaires, arborescents, déconnectés et avec cycles) ;
    \item \textbf{Tests Frontend :} Spécifications sous Jasmine et Karma validant le rendu et la gestion d'état des composants Angular ;
    \item \textbf{Validation Expérimentale sur Banc d'Essai :} Rejeu de scénarios d'incidents via le simulateur pour mesurer les temps de réponse de bout en bout.
\end{itemize}

\section{Suite de Tests Unitaires Backend (.NET 9 / xUnit)}

Le projet de test \texttt{DigitalTwin.Domain.Tests} intègre \textbf{83 tests unitaires automatisés} exécutés sous le framework \textbf{xUnit} avec les assertions expressives de \textbf{FluentAssertions}. Le Tableau~\ref{tab:tests_backend} synthétise la répartition des cas de test.

\begin{table}[htbp]
\centering
\small
\caption{Répertoire et répartition des 83 tests unitaires du backend}
\label{tab:tests_backend}
\begin{tabularx}{\textwidth}{l Y c c}
\toprule
\textbf{Composant / Classe Testée} & \textbf{Objectif de Validation} & \textbf{Nb Tests} & \textbf{Statut} \\
\midrule
\texttt{MachineTests} & Invariants d'agrégat, transitions d'état, règles de seuils & 22 & Validé (100\%) \\
\texttt{SensorTests} & Validation des unités, cohérence Warning/Critical & 18 & Validé (100\%) \\
\texttt{AnomalyDetectionServiceTests} & Qualification des plages de déviation (Nominal/Warning/Critical) & 15 & Validé (100\%) \\
\texttt{AlertCooldownServiceTests} & Fenêtre glissante anti-rebond, concurrence d'accès & 12 & Validé (100\%) \\
\texttt{ImpactAnalysisServiceTests} & Algorithme BFS, acyclicité, profondeur, topologies complexes & 16 & Validé (100\%) \\
\midrule
\textbf{Total} & \textbf{Couverture globale du cœur de domaine} & \textbf{83} & \textbf{Succès} \\
\bottomrule
\end{tabularx}
\end{table}

\section{Validation de l'Algorithme de Propagation BFS}

L'algorithme de calcul d'impact a été soumis à des bancs de tests validant ses propriétés fondamentales :
\begin{itemize}
    \item \textbf{Complétude :} Toutes les machines avals connectées à la source sont identifiées sans omission ;
    \item \textbf{Exactitude de la profondeur :} La distance topologique $d$ calculée correspond strictement au plus court chemin dans le graphe orienté ;
    \item \textbf{Résistance aux cycles :} En introduisant artificiellement une boucle de dépendance ($A \rightarrow B \rightarrow C \rightarrow A$), l'ensemble \texttt{visited} bloque la ré-insertion du sommet dans la file $Q$, évitant toute boucle infinie ou débordement de pile (\textit{StackOverflow}).
\end{itemize}

\section{Validation Expérimentale sur le Scénario Concasseur CR-001}

Le scénario de dégradation thermique et vibratoire du concasseur primaire \texttt{CR-001} a été exécuté sur le banc d'essai. La Figure~\ref{fig:scenario_cr001} illustre l'évolution temporelle de la mesure de vibration et le comportement dynamique du système de jumeau numérique.

\begin{figure}[htbp]
\centering
\begin{tikzpicture}
\begin{axis}[
    width=13.5cm,
    height=6.5cm,
    xlabel={Temps de simulation (secondes)},
    ylabel={Vibration (mm/s)},
    xmin=0, xmax=120,
    ymin=0, ymax=10,
    grid=both,
    grid style={line width=.1pt, draw=gray!20},
    major grid style={line width=.2pt, draw=gray!40},
    legend pos=north west,
    legend cell align={left},
    legend style={font=\footnotesize\sffamily}
]
    % Courbe de mesure de vibration simulée
    \addplot[
        color=ocpnavy,
        thick,
        mark=none,
        smooth
    ] coordinates {
        (0, 2.2) (10, 2.5) (20, 2.8) (30, 3.4) 
        (40, 4.2) (45, 4.8) (50, 5.6) (60, 6.8) 
        (65, 7.6) (70, 8.4) (80, 8.8) (90, 7.8) 
        (95, 6.2) (100, 4.6) (110, 3.0) (120, 2.3)
    };
    \addlegendentry{Mesure Capteur (CR-001)}

    % Seuil Warning (4.5 mm/s)
    \addplot[
        color=stwarning,
        dashed,
        very thick
    ] coordinates {(0, 4.5) (120, 4.5)};
    \addlegendentry{Seuil Warning ($4.5$\,mm/s)}

    % Seuil Critical (7.0 mm/s)
    \addplot[
        color=stcritical,
        dotted,
        very thick
    ] coordinates {(0, 7.0) (120, 7.0)};
    \addlegendentry{Seuil Critical ($7.0$\,mm/s)}

    % Annotations textuelles des phases
    \node[font=\tiny\bfseries, color=stwarning, anchor=south] at (axis cs:48, 5.0) {Alerte Warning (t=44s)};
    \node[font=\tiny\bfseries, color=stcritical, anchor=south] at (axis cs:68, 7.8) {Alerte Critique + BFS (t=62s)};
    \node[font=\tiny\bfseries, color=stnominal, anchor=south] at (axis cs:108, 3.2) {Retour Nominal (t=105s)};
\end{axis}
\end{tikzpicture}
\caption{Évolution temporelle de la vibration de CR-001 et franchissement dynamique des seuils}
\label{fig:scenario_cr001}
\end{figure}

Les observations relevées lors de l'expérimentation confirment la conformité fonctionnelle :
\begin{enumerate}
    \item À $t = 44$\,s, la vibration dépasse $4.5$\,mm/s. Une alerte \textit{Warning} est générée instantanément et l'état du concasseur bascule en \textit{Warning} sur la cartographie topologique ;
    \item À $t = 62$\,s, la vibration franchit le seuil critique de $7.0$\,mm/s. Le \texttt{DigitalTwinEngine} déclenche l'analyse BFS : le convoyeur \texttt{CV-001} ($d=1$) et le crible \texttt{SC-001} ($d=2$) passent immédiatement en état \textit{Degraded} sur l'interface opérateur ;
    \item Le mécanisme d'anti-rebond (\textit{Cooldown}) empêche l'émission de notifications superflues durant la phase de plateau critique ($t = 65$\,s à $85$\,s) ;
    \item À $t = 105$\,s, la vibration redescend en zone nominale. L'équipement se réinitialise et l'ensemble de la chaîne retrouve son état vert nominal.
\end{enumerate}

\section{Mesures de Performance et Latence Système}

Les temps de traitement ont été profilés sur 1\,000 requêtes consécutives de télémétrie :
\begin{itemize}
    \item \textbf{Temps d'exécution du pipeline backend (étapes 1 à 6) :} $12.4$\,ms en moyenne ;
    \item \textbf{Temps d'exécution du parcours BFS (graphe de 20 nœuds) :} $1.2$\,ms ;
    \item \textbf{Latence de diffusion WebSocket SignalR vers le client :} $32.0$\,ms ;
    \item \textbf{Temps de réponse total perçu par l'opérateur :} $45.6$\,ms (largement inférieur à l'exigence de $100$\,ms).
\end{itemize}

\section{Synthèse de Validation des Besoins}

Toutes les exigences fonctionnelles (BF-01 à BF-08) et non-fonctionnelles (BNF-01 à BNF-05) définies au cahier des charges sont intégralement couvertes et validées par les tests automatisés et expérimentaux.
